\chapter{投后管理与价值创造}

\section{从监控到主动价值创造}

投后不仅是收取报表：优秀 GP 在交易完成前已草拟\textbf{价值创造路线图}（采购协同、定价、渠道、数字化、组织瘦身等）。执行依赖被投企业管理层的信任与激励——投后首 100 天的沟通节奏与「不做什么」的承诺同样重要。

\section{财务与运营 KPI}

建立与战略一致的 KPI 仪表盘：收入增长、利润率、现金转换周期、客户留存、安全与环境指标（若适用）。对比投资备忘录中的基准，定期复盘偏差原因；必要时触发\textbf{预警会}与纠偏措施。

\section{人才与组织}

关键人风险是 PE 失败主因之一。投后应参与\textbf{高管评估、薪酬与长期激励}设计，必要时引入猎头与外部董事。文化冲突在跨境并购中尤为突出，需预留整合预算与时间。

\section{追加投资与跟投}

后续轮次是否跟进，应基于\textbf{重新估值}与\textbf{权利维持}（防稀释），而非沉没成本。LP 侧共同投资（co-invest）机会需公平分配并记录，避免利益输送指控。

\section{危机项目}

业绩持续恶化、治理僵局或合规爆雷时，应启动\textbf{危机协议}：增派资源、更换管理层、寻求出售或重组。法律上注意董事\textbf{受信义务}在困境中的变化。早期识别比晚期救火成本更低；投后例会应保留「黄色/红色」分级机制。
